%% ma: L11-B17-0040
%% lop: 11
%% chuong: 5
%% bai: 17
%% dang: 11.17.2
%% loai: TN4
%% muc: VD
%% dap_an: "D"
%% nguon: "tai-lieu-hoc-tap-mon-toan-lop-11-theo-chuong-trinh-giao-duc-pho-thong-2018.pdf (Lớp Toán Thầy Duy, Phạm Lê Duy, THPT Chu Văn An), Chương 3 Bài 3, Phần 1, Câu 20 (tr.405)"
%% kien_thuc: "Lớp 11 Bài 17 (hàm số liên tục)"
%% trang_thai: cho-duyet
%% ly_do: "Dạng mới đề xuất (11.17.x), chờ giáo viên duyệt dạng; tài liệu gốc không có lời giải mẫu, đáp án do kho tự giải và kiểm bằng sympy."
\begin{ex}
Cho hàm số $f(x)=\begin{cases}\dfrac{\sqrt{x+3}-m}{x-1}&\text{khi }x\ne1\\n&\text{khi }x=1\end{cases}$. Để hàm số liên tục tại $x_0=1$ thì giá trị của biểu thức $(m+n)$ tương ứng bằng
\choice{$\dfrac34$.}{$1$.}{$-\dfrac12$.}{\True $\dfrac94$.}
\loigiai{\sloppy Để giới hạn hữu hạn thì tử phải triệt tiêu tại $x=1$: $\sqrt4-m=0\Rightarrow m=2$. Khi đó $\dfrac{\sqrt{x+3}-2}{x-1}=\dfrac1{\sqrt{x+3}+2}\to\dfrac14$, nên $n=\dfrac14$. Vậy $m+n=\dfrac94$.}
\end{ex}
